\section{Agent 的本体：ReAct 循环、工具调用与测试时计算}

\subsection{从 buzzword 回到最小定义}
Sutton 在 00:04 附近直接给出自己的定义：所谓 LLM Agent，本质是\textbf{动态计算 + 工具使用}。动态计算意味着模型不是一次输出最终答案，而是按任务难度分配更多测试时计算；工具使用意味着它可以把中间假设外化到可执行环境，通过命令、代码运行、调试信息不断修正下一步策略。

\begin{knowledgebox}{最小可运行 Agent 循环}
\[
\text{Trajectory}_{t}
\xrightarrow{\text{LLM}}
\text{Thought + Tool Call}
\xrightarrow{\text{Environment}}
\text{Observation}_{t+1}
\xrightarrow{\text{append}}
\text{Trajectory}_{t+1}
\]
这个循环与 ReAct 思路一致：推理不是终点，而是为了决定下一次工具动作。
\end{knowledgebox}

\begin{importantbox}{为什么“动态计算”比“单次聪明”更关键}
\begin{itemize}
    \item 复杂任务往往前几步就会犯错，允许迭代才能暴露并纠正错误路径。
    \item 同一模型在 ``多轮试错 + 可执行反馈'' 下，能力上限显著高于 ``一次生成''。
    \item 对真实工程问题，持续验证比单次文本流畅更有价值。
\end{itemize}
\end{importantbox}

这套定义与课堂中反复出现的句子一致：\textit{``Chain of Thought plus tool use is at the heart of agents.''} 在 coding 场景里，工具输出就是下一轮推理的训练信号；在安全场景里，调试器、sanitizer、程序崩溃日志则成为更强的外部反馈。

\begin{figure}[H]
\centering
\includegraphics[width=0.88\textwidth]{frame-07-react-loop.jpg}
\caption{课程后段再次回到 ReAct 风格循环：观察、分析、调用工具、再观察 \protect\footnotemark}
\end{figure}
\footnotetext{视频画面时间区间：01:18:06--01:18:40。}

\subsection{测试时计算预算：继续修补还是重启轨迹}
讲座中一个很实用的问题是：当得到一个 ``部分正确但未通过'' 的 patch，应该继续在当前轨迹修补，还是丢弃后重新采样？这其实是在做\textbf{测试时计算预算分配}。对于短小改动，重新采样可能更便宜；对于复杂多文件变更，保留上下文继续修补可能更优。

\begin{warningbox}{常见误区：只看单条轨迹质量}
如果只分析最好那条轨迹，很容易高估系统。真实部署中应同时看：
\begin{itemize}
    \item 每次成功平均消耗的轨迹数；
    \item 失败轨迹是否可恢复；
    \item 超时或循环失败的比例。
\end{itemize}
\end{warningbox}

\subsection{工具不是附属功能，而是能力边界}
在课堂表述里，工具调用不是 ``给模型一点插件''，而是直接决定问题可解空间。没有代码浏览工具，模型几乎无法稳定定位大仓库问题；没有执行反馈工具，模型会长期停留在语义猜测；没有调试能力，安全场景中的漏洞触发链难以闭环。

\begin{importantbox}{工具调用的系统含义}
一个 Agent 的能力上限并不只由 base model 决定，还由 ``工具集合 $\times$ 反馈格式 $\times$ 控制策略'' 共同定义。主讲人把这三者称为需要一起调参的系统层面设计。
\end{importantbox}

\subsection{本章小结}
Agent 在本讲中的定义很明确：它是一个可持续迭代的动作系统，而非长回答生成器。动态计算、工具反馈、可恢复轨迹构成了后续所有工程讨论的基础。

